\documentclass[12pt]{article}

\usepackage[utf8]{inputenc}
\usepackage[spanish]{babel}
\usepackage{geometry}
\geometry{margin=2.5cm}

\title{UNIVERSIDAD TÉCNICA ESTATAL DE QUEVEDO 
\\FORM-FB-01 \\ Fichas Bibliográficas \\ Fundamentos de Sistemas Distribuidos y Arquitecturas Distribuidas}
\author{Farinango Guandinango Freddy }
\date{2026}

\begin{document}

\maketitle
\section*{Ficha Bibliográfica 1}

\textbf{Referencia APA 7}

Chen, H., Qi, W., Li, W., Shuai, S., Zheng, F., \& Yu, Z. (2021). Design and Implementation of Gray Publishing System under Distributed Application Microservice. \textit{2021 11th International Conference on Power and Energy Systems, ICPES 2021}, 452–455. 
https://doi.org/10.
1109/ICPES53652.2021.9683953

\vspace{0.3cm}

\textbf{DOI o URL}

https://doi.org/10.1109/ICPES53652.2021.9683953

\vspace{0.3cm}

\textbf{Resumen}

Este artículo presenta el diseño e implementación de un sistema de publicación basado en arquitectura de microservicios dentro de aplicaciones distribuidas. Los autores explican cómo el uso de microservicios mejora la escalabilidad, el mantenimiento y la flexibilidad del sistema. Además, se describe la separación de funciones en distintos servicios independientes, permitiendo una administración más eficiente y una mejor capacidad de actualización. El estudio demuestra que las arquitecturas distribuidas basadas en microservicios pueden optimizar el rendimiento y la estabilidad de aplicaciones modernas. También se destacan beneficios relacionados con la modularidad y la reducción de errores en sistemas complejos. Finalmente, los autores concluyen que este enfoque facilita el desarrollo y la integración de nuevas funcionalidades.

\vspace{0.3cm}

\textbf{Nota personal}

Gracias a este  artículo me permitió comprender cómo los microservicios facilitan la división de aplicaciones complejas en servicios independientes, mejorando la escalabilidad y el mantenimiento de sistemas distribuidos. 


\section*{Ficha Bibliográfica 2}

\textbf{Referencia APA 7}

Jiang, Y., Zhang, N., \& Ren, Z. (2020). Research on intelligent monitoring scheme for microservice application systems. \textit{Proceedings - 2020 International Conference on Intelligent Transportation, Big Data and Smart City, ICITBS 2020}, 791–794. https://doi.org/10.1109/ICITBS49701.2020.00173

\vspace{0.3cm}

\textbf{DOI o URL}

https://doi.org/10.1109/ICITBS49701.2020.00173  

\vspace{0.3cm}

\textbf{Resumen}

El artículo analiza un esquema inteligente de monitoreo para sistemas de aplicaciones basados en microservicios. Los autores explican la importancia de supervisar el rendimiento, la estabilidad y la comunicación entre servicios distribuidos. También describen técnicas para detectar errores y optimizar el funcionamiento del sistema en tiempo real. El estudio muestra cómo el monitoreo inteligente ayuda a mejorar la confiabilidad y disponibilidad de aplicaciones modernas. Además, se resalta que las arquitecturas de microservicios requieren mecanismos avanzados de supervisión debido a la complejidad de la interacción entre múltiples servicios independientes. Finalmente, el trabajo destaca que estas estrategias permiten identificar fallas rápidamente y reducir tiempos de respuesta en entornos distribuidos.

\vspace{0.3cm}

\textbf{Nota personal}

Mediante este artículo se logró identificar la importancia del monitoreo en sistemas distribuidos y la necesidad de herramientas de supervisión en arquitecturas de microservicios. 



\section*{Ficha Bibliográfica 3}

\textbf{Referencia APA 7}

Goudarzi, M., Wu, H., Palaniswami, M., \& Buyya, R. (2021). An Application Placement Technique for Concurrent IoT Applications in Edge and Fog Computing Environments. \textit{IEEE Transactions on Mobile Computing}, \textit{20}(4), 1298–1311. https://doi.org/10.1109/TMC.2020.29
67041

\vspace{0.3cm}

\textbf{DOI o URL}

https://doi.org/10.1109/TMC.2020.2967041

\vspace{0.3cm}

\textbf{Resumen}

Este artículo propone una técnica de ubicación de aplicaciones para entornos Edge y Fog Computing orientados a aplicaciones IoT concurrentes. Los autores estudian cómo distribuir aplicaciones de manera eficiente para reducir la latencia y mejorar el rendimiento. También se analizan problemas relacionados con la asignación de recursos y la optimización de tareas en sistemas distribuidos. El estudio demuestra que las técnicas inteligentes de distribución permiten mejorar el uso de recursos computacionales y la calidad del servicio. Además, se destaca la relevancia de Edge y Fog Computing en aplicaciones modernas conectadas al Internet de las Cosas. Finalmente, los autores concluyen que estas tecnologías favorecen el procesamiento rápido de datos y aumentan la eficiencia de sistemas distribuidos complejos.

 


\vspace{0.3cm}

\textbf{Nota personal}

El estudio investigado permitió comprender cómo las aplicaciones IoT pueden distribuirse eficientemente en entornos Edge y Fog Computing para optimizar el rendimiento y reducir la latencia. 


\section*{Ficha Bibliográfica 4}

\textbf{Referencia APA 7}

Goel, A., \& Thangaraju, B. (2022). Authenticating Distributed Systems Using SPIRE over Kubernetes Cluster. \textit{2022 IEEE International Conference on Electronics, Computing and Communication Technologies, CONECCT 2022}. https://doi.org/10.1109/CONECCT55679.2
022.9865835


\vspace{0.3cm}

\textbf{DOI o URL}

https://doi.org/10.1109/CONECCT55679.2022.9865835

\vspace{0.3cm}

\textbf{Resumen}

El artículo explica cómo utilizar SPIRE y SPIFFE para autenticar sistemas distribuidos dentro de clústeres Kubernetes. Los autores describen la importancia de implementar seguridad en arquitecturas de microservicios y entornos de nube. Además, se presentan mecanismos de autenticación mutua mediante certificados criptográficos para garantizar comunicaciones seguras entre servicios. El estudio demuestra cómo estas tecnologías ayudan a construir entornos Zero Trust y prevenir accesos no autorizados. También se analiza la integración de SPIRE en Kubernetes para mejorar la seguridad y confiabilidad de aplicaciones distribuidas modernas. Finalmente, el trabajo resalta que estas herramientas fortalecen la protección de datos y reducen riesgos de vulnerabilidades en aplicaciones distribuidas.

\vspace{0.3cm}

\textbf{Nota personal}

El artículo permitió analizar la importancia de la autenticación y seguridad en sistemas distribuidos basados en Kubernetes y arquitecturas de microservicios. 



\section*{Ficha Bibliográfica 5}

\textbf{Referencia APA 7}

Söylemez, M., Tekinerdogan, B., \& Tarhan, A. K. (2022). Challenges and Solution Directions of Microservice Architectures: A Systematic Literature Review. \textit{Applied Sciences 2022, Vol. 12, Page 5507}, \textit{12}(11), 5507. https://doi.org/10.3390/APP12115507

\vspace{0.3cm}

\textbf{DOI o URL}

https://doi.org/10.3390/app12115507

\vspace{0.3cm}

\textbf{Resumen}

Este artículo presenta una revisión sistemática de literatura sobre los desafíos y soluciones relacionados con arquitecturas de microservicios. Los autores identifican problemas asociados con escalabilidad, monitoreo, seguridad, comunicación y mantenimiento en sistemas distribuidos. También se analizan diferentes estrategias y soluciones propuestas en investigaciones recientes para enfrentar estos desafíos. El estudio demuestra que las arquitecturas de microservicios ofrecen beneficios importantes, pero requieren planificación y herramientas adecuadas para garantizar su correcto funcionamiento. Además, se destaca la importancia de DevOps, automatización y diseño orientado a dominios en el desarrollo de aplicaciones modernas. Finalmente, los autores concluyen que la adopción de buenas prácticas mejora la eficiencia, confiabilidad y adaptación de sistemas distribuidos complejos.

\vspace{0.3cm}

\textbf{Nota personal}

A través de esta investigación se identificaron los principales desafíos de las arquitecturas de microservicios y las estrategias utilizadas para mejorar su estabilidad y mantenimiento. 



\end{document}
